% 附录 A：高观点结论的考纲合法化转译手册
\section{附录 A：高观点结论的考纲合法化转译手册}\label{app:legal}

\begin{commonerror}{铁律：草稿纸用高观点，答题卡写初等通法}
解答题中\textbf{严禁}出现「由仿射变换知」「由极点极线知」「由欧拉公式知」「由 Abel 变换知」「由洛必达法则知」等超纲字眼。这类词一旦出现，轻则扣步骤分，重则该问判零。高观点的正确用法是：在草稿纸上用它锁定答案、取等位置与参数边界，然后在卷面上用考纲认可的初等推导把结论\textbf{重新证一遍}——代数验算本身就是满分证明。
\end{commonerror}

\begin{conccard}[三级]{三类角色各就各位}{阅卷老师只看卷面}{把预判当成证明}
\textbf{草稿纸}：随便用，越快越准越好；\textbf{卷面}：只允许考纲工具（联立、韦达定理、作差构造、错位相减、分类讨论、切线不等式）；\textbf{采分点}：看关键等式有没有被显式写出。三者不冲突，缺的只是翻译。
\end{conccard}

下面六个模板覆盖高中压轴的高频场景，每个都按「预判 $\to$ 卷面 $\to$ 采分点」三段给出。

\subsection*{模板一：圆锥曲线中的仿射压缩}

\textbf{草稿纸预判}：椭圆沿短轴方向压缩 $b/a$ 倍即得圆，面积同比例缩小，故内接三角形最大面积为 $\dfrac{3\sqrt3}{4}ab$；斜率乘积定值 $-\dfrac{b^2}{a^2}$。

\textbf{卷面写法}：设过定点的直线为 $y=kx+m$，代入 $\dfrac{x^2}{a^2}+\dfrac{y^2}{b^2}=1$，写出韦达定理与判别式条件；由弦长公式 $|AB|=\sqrt{1+k^2}\,|x_1-x_2|$ 与点到直线距离写出面积
\[
  S=\frac12|AB|\cdot d ,
\]
整理成关于 $k,m$ 的函数，用基本不等式或求导在预判的取等点处取得最值，并验证此时满足判别式大于零。

\textbf{采分点}：韦达定理两式、弦长表达式、距离表达式、「当且仅当……时取等」以及取等条件的检验。仿射结论只用来决定往哪条直线上凑。

\subsection*{模板二：极点与极线的定点定值}

\textbf{草稿纸预判}：由极点 $P(x_0,y_0)$ 直接写出极线方程 $\dfrac{x_0x}{a^2}+\dfrac{y_0y}{b^2}=1$，目标直线即刻锁定，不必再猜。

\textbf{卷面写法}：设两交点 $A(x_1,y_1),B(x_2,y_2)$，分别代入曲线方程并作差：
\[
  \frac{x_1^2-x_2^2}{a^2}+\frac{y_1^2-y_2^2}{b^2}=0
  \ \Longrightarrow\
  \frac{(x_1+x_2)(x_1-x_2)}{a^2}+\frac{(y_1+y_2)(y_1-y_2)}{b^2}=0 .
\]
再结合 $A,B$ 都在已知直线上，整理出 $\dfrac{y_1-y_2}{x_1-x_2}$ 与 $\dfrac{y_1+y_2}{x_1+x_2}$ 的关系，最后代入要证的对称式化简即可。

\textbf{采分点}：两点式作差这一步必须显式写出；之后全是恒等变形。

\subsection*{模板三：不动点与特征根求通项}

\textbf{草稿纸预判}：解 $x^*=px^*+q$ 或 $\lambda^2-p\lambda-q=0$，读出公比与两根。

\textbf{卷面写法}：由原递推式恒等变形直接写出
\[
  a_{n+1}-x^*=p(a_n-x^*) ,
\]
或两根情形下
\[
  \frac{a_{n+1}-\lambda_1}{\,a_{n+1}-\lambda_2\,}=k\,\frac{a_n-\lambda_1}{\,a_n-\lambda_2\,} .
\]
\textbf{完全不需要解释} $x^*$、$\lambda_i$ 是怎么来的——作为「已知常数」使用，再用首项确定待定系数，最后验证 $n=1$ 成立。

\textbf{采分点}：构造式本身、首项代入、$n=1$ 的检验。

\subsection*{模板四：Abel 变换与错位相减}

\textbf{草稿纸预判}：用分部求和快速算出 $S_n$ 的闭式。

\textbf{卷面写法}：规范列出
\[
  S_n=a_1+a_2+\cdots+a_n,\qquad
  qS_n=\phantom{a_1+{}}a_1q+\cdots+a_nq ,
\]
错开一位对齐后写「两式相减得」，直接默写化简结果，并把 $n=1$ 与 $n\to\infty$（若涉及极限）单独验证。

\textbf{采分点}：两行原始展开、相减动作、最终合并式、首项检验。

\subsection*{模板五：三角恒等变换与欧拉公式}

\textbf{草稿纸预判}：忘了三倍角公式时，展开 $(\cos\theta+i\sin\theta)^3$ 取实部，秒得 $\cos3\theta=4\cos^3\theta-3\cos\theta$。

\textbf{卷面写法}：从两角和公式出发现场推导：
\[
  \cos3\theta=\cos(2\theta+\theta)=\cos2\theta\cos\theta-\sin2\theta\sin\theta=\cdots,
\]
再降幂为单角函数，最后写出所需形式。

\textbf{采分点}：起始公式明确、每一步降幂正确、最终形式与题目要求一致（平方关系、次数）。

\subsection*{模板六：导数压轴的隐零点与切线放缩}

\textbf{草稿纸预判}：用切线 $e^x\ge x+1$、$\ln x\le x-1$ 预判参数临界值 $a_0$；必要时用端点效应估出 $a_0$。

\textbf{卷面写法}：按 $a\le a_0$ 与 $a>a_0$ 分类讨论。对每一类，先证切线不等式（以 $e^x-x-1$ 求最小值的方式），再代入原函数完成放缩链；隐零点处写成「设 $x_0$ 满足 $f'(x_0)=0$，则……」，用整体代换避免解出具体数值。

\textbf{采分点}：分类标准清晰、切线不等式的完整证明、隐零点的设而不求与代换、端点值的检验。

\subsection*{自查清单}

\begin{itemize}
  \item 全篇是否出现超纲名词？若有，一律替换为等价代数操作。
  \item 关键等式是否都被显式写出（而不是「显然」「易得」）？
  \item 取等条件、判别式正负、$n=1$ 首项是否逐一检验？
  \item 参数分类是否有重叠或遗漏？
\end{itemize}

\begin{conccard}[一级]{一句话记住}{预判给方向，通法给分数}{两者都要}
\key{草稿纸负责快，卷面负责稳}。高观点决定你往哪条路上走，考纲通法决定你能拿多少分——只有前者，会丢步骤分；只有后者，会算不完。
\end{conccard}
